% ============================================================
%  body.tex - Experiment B: pi-0.5 on LIBERO
%  Section-only content (no preamble). Imported by
%      reports/experiments/2026.10.05-2026.10.11-pi05-libero/pi05-libero.tex
%      report.tex
%  The importing document must load the `ego-report` style first.
% ============================================================

\section{Experiment B: \texorpdfstring{$\pi_{0.5}$}{pi-0.5} on LIBERO}
\label{sec:pi05}

\subsection{Overview and Scope}
\label{sec:pi05-overview}

This experiment reproduces the official four-suite LIBERO success rate of
Physical Intelligence's \emph{$\pi_{0.5}$} policy~\cite{pi05_2025} on a single
A100 GPU. The authors' released LIBERO fine-tuned weights are evaluated with
the OpenPI reference implementation; no new $\pi_{0.5}$ model was trained, no
RoboDojo-specific weights were substituted, and no video-classification score
is used in place of a robot task success rate.

The main protocol, labelled P12, uses the authors' default JAX precision. A
separate control run P10 uses the highest precision with TF32 disabled and
reaches $97.00\,\%$; it is reported as a stability and precision control only
and does not replace the pre-registered main protocol. The P labels are internal
stage tags, not model versions: \emph{rpc} denotes a separated
model/simulation process, \emph{native} the authors' default protocol and
\emph{full} a 2000-episode run.

The evaluation covers the four standard suites --- Spatial, Object, Goal and
LIBERO-10 --- each with 10 tasks and 50 episodes per task, for
$40 \times 50 = 2000$ closed-loop episodes in total. The complete statistics,
the per-task breakdown and representative teaching cases are described below.

\subsection{Benchmark and Evaluation Protocol}
\label{sec:pi05-benchmark}

LIBERO is a language-conditioned manipulation benchmark for simulated robot
arms~\cite{libero2023}. Each task defines a tabletop scene, a natural-language
goal and a BDDL goal predicate; the policy observes the current RGB frames and
the robot state, predicts an action, and MuJoCo advances the scene to produce
the next observation until the goal is met or the budget is exhausted. This is a
closed loop in which the robot's own actions change the next frame --- it is not
one-shot captioning of a recorded video.

\begin{table}[H]
  \centering
  \small
  \caption{The four LIBERO suites and the official $\pi_{0.5}$ scores used as
  the comparison target. Official numbers come from the OpenPI release
  table~\cite{pi05_2025}, not from this run.}
  \label{tab:pi05-suites}
  \begin{tabular}{@{}llrrrr@{}}
    \toprule
    Suite & Focus & Tasks & Per task & Max steps & Official \\
    \midrule
    Spatial   & spatial relations in language   & 10 & 50 & 220 & 98.8\,\% \\
    Object    & different target objects        & 10 & 50 & 280 & 98.2\,\% \\
    Goal      & different goal states           & 10 & 50 & 300 & 98.0\,\% \\
    LIBERO-10 & longer-horizon combined tasks   & 10 & 50 & 520 & 92.4\,\% \\
    \midrule
    Total / mean & four suites, equal weight    & 40 & 2000 total & --- & 96.85\,\% \\
    \bottomrule
  \end{tabular}
\end{table}

The scalar values used in this report have precise meanings. The final LIBERO
action is a 3-D translation command, a 3-D rotation command and a 1-D gripper
command; it is not a set of joint angles and the numbers cannot be read as
metres. At 20\,Hz each control step is about $0.05$\,s of simulated time, so a
220-step budget is roughly 11 seconds of simulated time and is unrelated to
wall-clock runtime. Success rate is successful episodes over all episodes; all
failures are retained and no run is presented as complete unless every episode
finished.

The state vector is an end-effector position (3), an axis-angle orientation (3)
and two gripper joint values (2), for a total of 8 dimensions. The axis-angle
vector encodes the rotation axis in its direction and the angle in its length;
it is not a set of Euler angles. A positive or negative gripper value is fixed
by the joint-axis convention. As explained in Section~\ref{sec:pi05-model},
this particular LIBERO configuration does not feed the state into the network,
so the interface field must not be mistaken for a model input.

\subsection{Model and Configuration}
\label{sec:pi05-model}

$\pi_{0.5}$ is a vision-language-action model: its interface receives a language
task, two camera RGB streams and the robot state, and emits a continuous action
chunk. It is not treated here as a world model, because this evaluation produces
actions and generates no future RGB. The official configuration used is
\texttt{Pi0Config(pi05=True,} \texttt{action\_horizon=10,}
\texttt{discrete\_state\_input=False)}; the concrete module implementation lives
in the pinned files \btt{src/openpi/models/pi0.py} and
\btt{pi0\_config.py}, and no network is invented beyond them.

Source review clarifies an important nuance. \texttt{embed\_prefix} consumes the
image and task tokens, while \texttt{embed\_suffix} creates a state token only
when \texttt{not self.pi05}. Because this configuration additionally sets
\texttt{discrete\_state\_input=False}, the discrete state text is disabled as
well, so this LIBERO weight does not use the state as a neural condition. The
8-D state is still transformed by the interface, stored and used for batch
shape inference. Receiving a state at the interface is therefore not the same as
using it, and this conclusion must not be generalised to other $\pi_{0.5}$
weights whose default configuration includes a discrete state.

The fixed implementation uses a SigLIP \texttt{So400m/14} vision encoder, a
\texttt{gemma\_2b} language backbone and a \texttt{gemma\_300m} action expert;
these names are source configuration labels and are not a full parameter
inventory. Vision tokens and the task text form the prefix; the action expert
maps noise to continuous actions and injects the flow-matching time through
adaRMSNorm. The internal action dimension is 32 with a horizon of 10, and the
final adapter takes the first 7 dimensions. The maximum token length is 200 and
the missing third camera is masked out; only two real cameras are present. Each
query starts from Gaussian noise and is denoised by the authors' default 10
flow-matching updates; the first 5 steps of the chunk are executed before
replanning.

Intuitively, the model predicts \emph{how to move next}, not \emph{what the
image becomes after the move}. Training such a VLA therefore needs paired
camera, instruction and continuous robot action data; action categories from
human first-person video are not a substitute for 7-D arm control targets. This
work reproduces inference and evaluation only and does not retrain the backbone
or the action expert.

\subsection{Closed-Loop Data Pipeline}
\label{sec:pi05-pipeline}

\begin{enumerate}
  \item The tasks are read in a fixed suite order from the pinned LIBERO BDDL
        and language definitions. Task descriptions are never reworded to suit
        the model.
  \item After \texttt{env.reset()} the author-provided initial state with index
        0--49 is applied; the environment seed is fixed at 7.
  \item Ten steps of \texttt{[0,0,0,0,0,0,-1]} are executed so that objects
        settle. These ten steps are not charged against the policy budget.
  \item $256 \times 256$ RGB frames are read from \texttt{agentview\_image} and
        \texttt{robot0\_eye\_in\_hand\_image} and flipped along both axes
        (equivalent to a $180^{\circ}$ rotation), matching the authors'
        training preprocessing.
  \item The frames are aspect-preservingly padded to $224 \times 224$ as
        \texttt{uint8}. Both real cameras are kept; the missing third camera is
        zero-filled and masked in the official adapter.
  \item The end-effector position (3), the quaternion converted to axis-angle
        (3) and the gripper joints (2) form the 8-D state. It is genuine robot
        proprioception; this configuration does not encode it as a network
        condition but records it for the interface and the audit.
  \item Quantile normalisation uses the statistics shipped with the checkpoint,
        the tokenizer encodes the task text, and the state is padded to the
        model's internal dimension.
  \item A $10 \times 7$ action chunk is predicted and checked for shape and
        finite values. The de-normalised output is executed directly, with no
        extra gripper inversion and no action format borrowed from another
        robot.
  \item Observations and predictions are refreshed every 5 steps. All predicted
        chunks, executed actions, states, rewards, \texttt{done} flags and the
        independent goal check are retained.
  \item The official \texttt{done} signal is the primary success criterion and
        \texttt{env.check\_success()} is a step-wise diagnostic. Reaching the
        maximum number of steps without success is recorded as a failure and is
        never removed from the denominator.
\end{enumerate}

A single query produces 10 steps but only the first 5 are executed before
replanning. PNG files are therefore indexed by query, while the video and
\texttt{steps.json} are indexed by the actual control step; the 5 unexecuted
steps remain predictions and must not be read as executed actions. The action
components $a_0$--$a_5$ are controller commands and $a_6$ is the gripper
command, so they cannot be interpreted directly as end-effector metres.

Figure~\ref{fig:pi05-dataflow} summarises the runtime flow.

\begin{figure}[H]
  \centering
  \fbox{\parbox{0.92\linewidth}{\centering\small
    two current cameras (256 RGB) $\rightarrow$ $180^{\circ}$ rotation and
    224 padding;\\
    EEF position / orientation + gripper 8-D state $\rightarrow$ quantile
    normalisation and padding;\\
    task language $\rightarrow$ tokenizer;\\
    (images, text) $\rightarrow$ vision-language backbone $\rightarrow$ action
    expert with 10 flow-matching steps;\\
    $\rightarrow$ $10 \times 7$ action chunk $\rightarrow$ execute first 5
    $\rightarrow$ new MuJoCo observation and \texttt{done}.}}
  \caption{Runtime data flow of the $\pi_{0.5}$ closed loop. The 8-D state is
  retained by the interface for the audit but is not a network condition in
  this configuration.}
  \label{fig:pi05-dataflow}
\end{figure}

\subsection{Frozen Source and Environment}
\label{sec:pi05-env}

\begin{table}[H]
  \centering
  \small
  \caption{Frozen source and environment of the main (P12) protocol.}
  \label{tab:pi05-env}
  \begin{tabular}{@{}p{0.32\linewidth}p{0.62\linewidth}@{}}
    \toprule
    Item & Value \\
    \midrule
    OpenPI revision & \btt{215abfb217dbac7d5f1273282331b9b1866c0479} \\
    LIBERO submodule & \btt{f78abd68ee283de9f9be3c8f7e2a9ad60246e95c} \\
    Runtime & JAX implementation, parameters restored as \texttt{bfloat16} \\
    Python environment & Python 3.11, dependencies in isolated per-run trees \\
    MuJoCo / robosuite & 3.2.3 / 1.4.1 \\
    Rendering & A100 via EGL (dispatcher unpacked into the project, loaded per subprocess) \\
    Environment seed & 7 \\
    Policy seed & 0, one independent policy per suite \\
    Initial states & the authors' fixed states 0--49 per task \\
    Action execution & predict 10, execute first 5, then replan \\
    \bottomrule
  \end{tabular}
\end{table}

The server reuses an existing Python 3.11 environment with dependencies
installed under isolated \texttt{runtime/} trees, without modifying the base
environment. MuJoCo 3.2.3 and robosuite 1.4.1 match the official lock file
consulted. Python, Torch and NumPy versions differ from the authors' older
client environment, so bitwise identity is not claimed.

EGL lets the A100 run MuJoCo without IsaacSim's RT-core requirement. The server
originally lacked the EGL dispatcher, so Ubuntu's official
\texttt{libegl1}/\texttt{libopengl0} packages were unpacked into the project's
\texttt{system-libs} directory and loaded only through a subprocess
\texttt{LD\_LIBRARY\_PATH}; no system package was installed and no driver was
changed. Version 2.3.7 was used only for a rendering probe; the policy cases and
the full evaluation all used 3.2.3.

The main-protocol run IDs are
\btt{pi05-p12-rpc-native-full-\{spatial,object,goal,10\}-s7}. The model and
simulation processes are launched through a project launcher that sets up the
isolated dependencies and the EGL environment; the launcher refuses to reuse an
existing run directory, so re-running requires a fresh unique run ID, a free
port and a new space/ETA registration. The command form is:

\begin{verbatim}
python launch_rpc_run.py \
  --config <runtime>/execution/<git>-pi05-native/spatial-config.json
\end{verbatim}

Each run records both subprocess commands and the explicit environment
overrides in \btt{command.json}, the frozen configuration, source
fingerprints, seed and budget in \btt{prereg.json}, and the exit code and
total wall time in \btt{status.json}.

\subsection{Results}
\label{sec:pi05-results}

The main P12 protocol completed all 2000 episodes and was independently
verified: $1935/2000 = 96.75\,\%$ versus the official $96.85\,\%$, a difference
of $-0.10$ percentage points. All 64{,}761 inferences and 320{,}142 executed
action steps were checked for finite values, with no invalid prediction.

\begin{table}[H]
  \centering
  \small
  \caption{Per-suite results. P12 is the pre-registered main protocol at the
  authors' default precision; P10 is a separate highest-precision control and
  does not replace the main table.}
  \label{tab:pi05-scores}
  \begin{tabular}{@{}lrrrrr@{}}
    \toprule
    Suite & Success / total & P12 & Official & Difference (pp) & P10 control \\
    \midrule
    spatial   & 495 / 500 & 99.00\,\% & 98.80\,\% & $+0.20$ & 97.60\,\% \\
    object    & 490 / 500 & 98.00\,\% & 98.20\,\% & $-0.20$ & 98.60\,\% \\
    goal      & 491 / 500 & 98.20\,\% & 98.00\,\% & $+0.20$ & 98.00\,\% \\
    LIBERO-10 & 459 / 500 & 91.80\,\% & 92.40\,\% & $-0.60$ & 93.80\,\% \\
    \midrule
    \textbf{Mean} & \textbf{1935 / 2000} & \textbf{96.75\,\%} &
    \textbf{96.85\,\%} & $\mathbf{-0.10}$ & \textbf{97.00\,\%} \\
    \bottomrule
  \end{tabular}
\end{table}

A precision control on the same Spatial inputs with policy seed 0 confirmed that
the authors' default and the historical highest/TF32-off settings are not
numerically identical: the maximum action difference was $0.007402592285156384$.
Because the official entry point does not set a JAX matrix precision, P12 keeps
the authors' default precision and executes revision \texttt{c667a73}.

Figure~\ref{fig:pi05-official} compares the official, main-protocol and
precision-control scores, and Figure~\ref{fig:pi05-40tasks} shows the per-task
rates.

\begin{figure}[H]
  \centering
  \includegraphics[width=0.85\linewidth]{2026.10.05-2026.10.11-pi05-libero/official-control-native.png}
  \caption{Official, main-protocol (P12) and precision-control (P10) scores.}
  \label{fig:pi05-official}
\end{figure}

\begin{figure}[H]
  \centering
  \includegraphics[width=0.95\linewidth]{2026.10.05-2026.10.11-pi05-libero/all-40-task-rates.png}
  \caption{Per-task success rate over all 40 LIBERO tasks.}
  \label{fig:pi05-40tasks}
\end{figure}

{\footnotesize
\begin{longtable}{@{}l r p{6.6cm} r r >{\raggedright\arraybackslash}p{1.6cm}@{}}
  \caption{Per-task results (50 episodes per row). All initial states are
  retained; the failing IDs are kept for auditing and must not be used to select
  inputs for a future evaluation.}
  \label{tab:pi05-40tasks}\\
  \toprule
  Suite & ID & Task language & P12 /50 & P10 /50 & P12 failing IDs \\
  \midrule
  \endfirsthead
  \multicolumn{6}{@{}l}{\small\itshape Table~\ref{tab:pi05-40tasks}, continued}\\
  \toprule
  Suite & ID & Task language & P12 /50 & P10 /50 & P12 failing IDs \\
  \midrule
  \endhead
  \midrule
  \multicolumn{6}{r@{}}{\small\itshape continued on the next page}\\
  \endfoot
  \bottomrule
  \endlastfoot
    spatial & 0 & pick up the black bowl between the plate and the ramekin and place it on the plate & 50 & 50 & --- \\
    spatial & 1 & pick up the black bowl next to the ramekin and place it on the plate & 50 & 50 & --- \\
    spatial & 2 & pick up the black bowl from table center and place it on the plate & 50 & 50 & --- \\
    spatial & 3 & pick up the black bowl on the cookie box and place it on the plate & 49 & 49 & 7 \\
    spatial & 4 & pick up the black bowl in the top drawer of the wooden cabinet and place it on the plate & 47 & 47 & 23, 25, 46 \\
    spatial & 5 & pick up the black bowl on the ramekin and place it on the plate & 50 & 47 & --- \\
    spatial & 6 & pick up the black bowl next to the cookie box and place it on the plate & 50 & 50 & --- \\
    spatial & 7 & pick up the black bowl on the stove and place it on the plate & 50 & 48 & --- \\
    spatial & 8 & pick up the black bowl next to the plate and place it on the plate & 50 & 50 & --- \\
    spatial & 9 & pick up the black bowl on the wooden cabinet and place it on the plate & 49 & 47 & 47 \\
    object & 0 & pick up the alphabet soup and place it in the basket & 48 & 49 & 2, 44 \\
    object & 1 & pick up the cream cheese and place it in the basket & 49 & 50 & 5 \\
    object & 2 & pick up the salad dressing and place it in the basket & 49 & 48 & 38 \\
    object & 3 & pick up the bbq sauce and place it in the basket & 49 & 49 & 33 \\
    object & 4 & pick up the ketchup and place it in the basket & 49 & 48 & 7 \\
    object & 5 & pick up the tomato sauce and place it in the basket & 48 & 50 & 7, 34 \\
    object & 6 & pick up the butter and place it in the basket & 50 & 50 & --- \\
    object & 7 & pick up the milk and place it in the basket & 49 & 50 & 1 \\
    object & 8 & pick up the chocolate pudding and place it in the basket & 49 & 50 & 16 \\
    object & 9 & pick up the orange juice and place it in the basket & 50 & 49 & --- \\
    goal & 0 & open the middle drawer of the cabinet & 47 & 49 & 1, 19, 31 \\
    goal & 1 & put the bowl on the stove & 50 & 50 & --- \\
    goal & 2 & put the wine bottle on top of the cabinet & 50 & 50 & --- \\
    goal & 3 & open the top drawer and put the bowl inside & 46 & 44 & 4, 20, 38, 45 \\
    goal & 4 & put the bowl on top of the cabinet & 50 & 50 & --- \\
    goal & 5 & push the plate to the front of the stove & 49 & 50 & 44 \\
    goal & 6 & put the cream cheese in the bowl & 49 & 49 & 49 \\
    goal & 7 & turn on the stove & 50 & 50 & --- \\
    goal & 8 & put the bowl on the plate & 50 & 50 & --- \\
    goal & 9 & put the wine bottle on the rack & 50 & 48 & --- \\
    libero\_10 & 0 & put both the alphabet soup and the tomato sauce in the basket & 48 & 46 & 18, 49 \\
    libero\_10 & 1 & put both the cream cheese box and the butter in the basket & 49 & 50 & 33 \\
    libero\_10 & 2 & turn on the stove and put the moka pot on it & 47 & 50 & 6, 25, 33 \\
    libero\_10 & 3 & put the black bowl in the bottom drawer of the cabinet and close it & 48 & 49 & 2, 13 \\
    libero\_10 & 4 & put the white mug on the left plate and put the yellow and white mug on the right plate & 46 & 49 & 7, 13, 34, 35 \\
    libero\_10 & 5 & pick up the book and place it in the back compartment of the caddy & 50 & 50 & --- \\
    libero\_10 & 6 & put the white mug on the plate and put the chocolate pudding to the right of the plate & 45 & 49 & 27, 30, 33, 38, 49 \\
    libero\_10 & 7 & put both the alphabet soup and the cream cheese box in the basket & 50 & 49 & --- \\
    libero\_10 & 8 & put both moka pots on the stove & 30 & 31 & 0, 2, 5, 6, 7, 9, 13, 14, 20, 21, 23, 26, 28, 30, 36, 37, 38, 41, 42, 49 \\
    libero\_10 & 9 & put the yellow and white mug in the microwave and close it & 46 & 46 & 3, 18, 44, 48 \\
\end{longtable}}

Acceptance checking compared, on the server, that no task or initial state was
missing or duplicated, that every episode's success agreed with the
\texttt{done}/\texttt{goal} signal, and that queries and actions matched the
stored NPZ files. The complete inputs, states and videos stay on the server; in
line with the delivery scope, only statistics, records and representative cases
were returned locally. Seeds were not selected and runs were not retried until
the score looked good, and agreement with the official number was not used as a
tuning target. The official release publishes only summary scores and no
per-episode predictions, so per-sample identity with the authors cannot be
claimed.

\subsection{Compute, Timing and Archive}
\label{sec:pi05-compute}

P12 ran from the first suite's start to the last suite's exit in 9673.44 seconds
(about 2\,h\,41\,m\,13\,s). The four suites overlap, so their wall times must
not be summed. This interval includes model loading, first compilation,
rendering, file writing and concurrency waiting, but excludes earlier downloads,
diagnostics, P10 and document preparation.

\begin{table}[H]
  \centering
  \small
  \caption{Per-suite runtime and request statistics. RPC latency includes
  networking, serialisation, service queuing and actual inference, so it is not
  a pure GPU kernel time.}
  \label{tab:pi05-timing}
  \begin{tabular}{@{}lrrrr@{}}
    \toprule
    Suite & Wall time & Queries & Control steps & RPC median / P95 (s) \\
    \midrule
    spatial   & 85.36 min  & 10{,}706 & 52{,}625  & 0.197 / 0.249 \\
    object    & 95.89 min  & 14{,}381 & 70{,}938  & 0.193 / 0.263 \\
    goal      & 85.43 min  & 11{,}660 & 57{,}342  & 0.211 / 0.265 \\
    LIBERO-10 & 147.45 min & 28{,}014 & 139{,}237 & 0.094 / 0.254 \\
    \midrule
    \textbf{Total} & \textbf{9673.44\,s} & \textbf{64{,}761} &
    \textbf{320{,}142} & --- \\
    \bottomrule
  \end{tabular}
\end{table}

At launch, five model/render process groups used roughly 46\,GiB in total and
then wound down as tasks finished. GPU sampling is whole-device, so a concurrent
peak must not be reported as a single-model memory footprint. The lower latency
in the second half of LIBERO-10 reflects reduced contention, not faster model
computation.

Each run directory stores \btt{summary.json} (500 episodes),
\btt{server/inference-audit.jsonl} (the finite-value check of every real model
call), \btt{prereg.json}, \btt{command.json}, \btt{status.json} and
\btt{gpu-samples.jsonl}. The returned local media contain, per episode:

\begin{table}[H]
  \centering
  \small
  \caption{Per-episode files and their boundaries.}
  \label{tab:pi05-files}
  \begin{tabular}{@{}p{0.30\linewidth}p{0.62\linewidth}@{}}
    \toprule
    File & Content and boundary \\
    \midrule
    \texttt{initial-state.npy} & the author initial state, used only for reset \\
    \texttt{query-NNN-*.png} & the two $224$ RGB inputs of one real query; not one PNG per control step \\
    \texttt{query-NNN.npz} & the 8-D state and the predicted $10 \times 7$ action; the last 5 steps may be unexecuted \\
    \texttt{queries.json} & every query index, control step, RPC wall latency and full prediction \\
    \texttt{steps.json} & the pre-step state, executed action and post-step reward/ \texttt{done}/goal \\
    \texttt{rollout.mp4} & the two RGB streams side by side, lossy; no extra terminal frame \\
    \texttt{summary.json} & task language, outcome, step and query counts, input fingerprints \\
    \bottomrule
  \end{tabular}
\end{table}

The action control rate is 20\,Hz in simulation while videos are written at
10\,fps for viewing, so a video lasts about twice the simulated action time.
Wall time additionally includes model, rendering, file-writing and concurrency
overhead; the three times must not be conflated. An archive is considered
complete only when both \texttt{archive-receipt.json} (private logs and
summaries) and \texttt{media/media-receipt.json} (per-file sizes and SHAs) exist
and are verified.

\subsection{Teaching Cases: Six Genuine Episodes}
\label{sec:pi05-cases}

The six P9 engineering cases below share a random stream across suites and are
not P12 main-table episodes; 5/6 success cannot stand in for the 2000-episode
score. First- and last-query images are not post-termination ground truth:
success is decided by the official \texttt{done}/\texttt{goal} signal after the
action. The complete \texttt{queries.json} stores all predictions, including
the unexecuted ones, while \texttt{steps.json} stores what was actually
executed.

Figure~\ref{fig:pi05-case-grid} collects the first genuine model input (external
and wrist view) of the spatial, object and goal cases, and
Figure~\ref{fig:pi05-case-actions} shows the executed actions with the observed
end-effector position for each case.

\begin{figure}[H]
  \centering
  \begin{subfigure}[t]{0.235\linewidth}
    \centering
    \includegraphics[width=\linewidth]{2026.10.05-2026.10.11-pi05-libero/spatial-t0-i0-first-agent.png}
    \caption{spatial t0 i0}
  \end{subfigure}\hfill
  \begin{subfigure}[t]{0.235\linewidth}
    \centering
    \includegraphics[width=\linewidth]{2026.10.05-2026.10.11-pi05-libero/object-t0-i0-first-agent.png}
    \caption{object t0 i0}
  \end{subfigure}\hfill
  \begin{subfigure}[t]{0.235\linewidth}
    \centering
    \includegraphics[width=\linewidth]{2026.10.05-2026.10.11-pi05-libero/goal-t0-i0-first-agent.png}
    \caption{goal t0 i0}
  \end{subfigure}\hfill
  \begin{subfigure}[t]{0.235\linewidth}
    \centering
    \includegraphics[width=\linewidth]{2026.10.05-2026.10.11-pi05-libero/goal-t0-i1-first-agent.png}
    \caption{goal t0 i1}
  \end{subfigure}
  \caption{First genuine external-camera input of four of the six teaching
  cases. The wrist views of the same queries are provided in the accompanying
  media folder.}
  \label{fig:pi05-case-grid}
\end{figure}

\begin{figure}[H]
  \centering
  \begin{subfigure}[t]{0.48\linewidth}
    \centering
    \includegraphics[width=\linewidth]{2026.10.05-2026.10.11-pi05-libero/spatial-t0-i0-actions.png}
    \caption{spatial t0 i0 (success, 101 steps)}
  \end{subfigure}\hfill
  \begin{subfigure}[t]{0.48\linewidth}
    \centering
    \includegraphics[width=\linewidth]{2026.10.05-2026.10.11-pi05-libero/spatial-t0-i1-actions.png}
    \caption{spatial t0 i1 (success, 76 steps)}
  \end{subfigure}

  \vspace{0.6em}
  \begin{subfigure}[t]{0.48\linewidth}
    \centering
    \includegraphics[width=\linewidth]{2026.10.05-2026.10.11-pi05-libero/object-t0-i0-actions.png}
    \caption{object t0 i0 (success, 131 steps)}
  \end{subfigure}\hfill
  \begin{subfigure}[t]{0.48\linewidth}
    \centering
    \includegraphics[width=\linewidth]{2026.10.05-2026.10.11-pi05-libero/object-t0-i1-actions.png}
    \caption{object t0 i1 (success, 146 steps)}
  \end{subfigure}

  \vspace{0.6em}
  \begin{subfigure}[t]{0.48\linewidth}
    \centering
    \includegraphics[width=\linewidth]{2026.10.05-2026.10.11-pi05-libero/goal-t0-i0-actions.png}
    \caption{goal t0 i0 (success, 116 steps)}
  \end{subfigure}\hfill
  \begin{subfigure}[t]{0.48\linewidth}
    \centering
    \includegraphics[width=\linewidth]{2026.10.05-2026.10.11-pi05-libero/goal-t0-i1-actions.png}
    \caption{goal t0 i1 (timeout failure, 300 steps)}
  \end{subfigure}
  \caption{Executed actions and observed end-effector position for the six
  teaching cases.}
  \label{fig:pi05-case-actions}
\end{figure}

The first genuine query of spatial task 0, initial state 0 is reproduced
verbatim below. It is real data, not a hypothetical example.

\begin{verbatim}
{
  "query_id": 0,
  "step": 0,
  "state_8d": [
    -0.21124162683394807, -0.011090601889321502, 1.1742495176982402,
    3.140895396142254, -0.002847819833825949, -0.08773931027790592,
    0.03872336505859045, -0.0387222020954563
  ],
  "first_predicted_action_7d": [
    0.010095496548414129, 0.027211936503648726, 0.05642917198252673,
    -0.01643640754377028, 0.005886954390445509, 0.0053137199111878974,
    -0.9974632556626797
  ]
}
\end{verbatim}

The position components are expressed in the simulation frame and are not
heights above the table; the gripper sign follows the joint-axis convention.
The state describes the present and the actions indicate the next step. Because
this configuration does not use the state as a network condition, the report
must not claim that the model consumed the 8-D state.

The six rollouts are stored alongside the figures as
\texttt{<suite>-t0-i<i>-rollout.mp4}. The complete predictions and executed
steps are provided as \texttt{queries.json} and \texttt{steps.json},
respectively. All inputs, predictions and videos were returned with per-file
SHA checks (\texttt{media-receipt.json}) and the six action-curve PNGs
(1{,}038{,}710 bytes in total) were verified as well. The action expert cannot be
credited with general robot generalisation or a paper-level contribution on the
basis of these six cases.

\subsection{Numerical Anomalies and Non-Finite Failures}
\label{sec:pi05-nan}

An early Object run stopped after the first 14 episodes of the P4 protocol when
a non-finite action appeared at initial state 14, exiting with code 1; the
partial run cannot be reported as a score and was retained as-is. A revision
then added diagnostic saving, counted an invalid prediction as an episode
failure and continued, without changing the gripper or the seed. The independent
500-episode retry was not spliced with the best results, and the supervisor
process started a retry only once, for the actually failing suite.

A long-horizon diagnostic read the GPU input tensors after 2699 anomalous
inferences and captured a device-side data anomaly: all three image tensors were
NaN, the tokens were abnormally zeroed and the raw 320 action output was NaN,
while the random noise was normal. Repeating with a fresh random key on the same
input still failed. The fact of the anomalous GPU input is confirmed, but the
underlying CUDA/EGL cause is not proven and must not be asserted. These
non-finite failures are distinct from genuine task failures under normal
numerical conditions.

\subsection{Normal Failures: Open Drawer, Still a Failure}
\label{sec:pi05-drawer}

The following two genuine episodes come from the P10 highest-precision control,
not from the P12 main table. Both use the language ``open the middle drawer of
the cabinet'', and the author BDDL goal
\texttt{Open wooden\_cabinet\_1\_middle\_region} requires the middle-drawer joint
coordinate to be below $-0.14$.

The original actions of 127 and 300 steps were replayed on the CPU without
rendering, in the original within-task environment reset order. For both
episodes the maximum difference over all 8-D end-effector states was 0, and the
step-wise \texttt{done} and \texttt{goal} signals agreed. The joint values below
therefore come from a verified physical replay.

\begin{table}[H]
  \centering
  \small
  \caption{Physical replay of the two drawer episodes.}
  \label{tab:pi05-drawer}
  \begin{tabular}{@{}lrrl@{}}
    \toprule
    Case & Middle drawer & Bottom drawer & Official goal \\
    \midrule
    task 0 / init 0 & $-0.14426769$ & $0.0$ & success \\
    task 0 / init 1 & $0.00147690$ & $-0.15999917$ & failure (bottom drawer opened) \\
    \bottomrule
  \end{tabular}
\end{table}

\begin{figure}[H]
  \centering
  \begin{subfigure}[t]{0.48\linewidth}
    \centering
    \includegraphics[width=\linewidth]{2026.10.05-2026.10.11-pi05-libero/drawer-success-agent.png}
    \caption{successful episode, last query}
  \end{subfigure}\hfill
  \begin{subfigure}[t]{0.48\linewidth}
    \centering
    \includegraphics[width=\linewidth]{2026.10.05-2026.10.11-pi05-libero/drawer-failure-agent.png}
    \caption{failed episode, last query}
  \end{subfigure}
  \caption{The last genuine query of the successful and the failing drawer
  episode. A visible open drawer does not prove task success.}
  \label{fig:pi05-drawer}
\end{figure}

A drawer is clearly open in both images, yet only one episode is successful.
Object identity and goal state must both be satisfied, so a picture of an open
drawer is not sufficient evidence of success. Every action in the failed
episode is finite: it is a genuine task failure and must not be grouped with the
earlier non-finite engineering anomalies.

The replay itself had failures worth recording. The first attempt used a wrong
import path. The first retry created a new environment per episode and produced
a maximum state error of $0.05557$ on the failing episode, so it cannot be
called an exact replay. The second retry preserved the official within-task
environment creation and reset order and reduced the error to zero. Random
sequence and internal controller state were not separated factor by factor, so
no specific internal causal mechanism is asserted. All three records were
returned with SHA and server acknowledgement and none was overwritten.

\subsection{Limitations and Next Steps}
\label{sec:pi05-limitations}

The claim of this experiment is bounded accordingly. Evaluation is closed-loop
action success in simulation only and involves no real-robot deployment. The
official numbers are summary scores with no per-episode archive, so per-sample
identity with the authors cannot be verified. The precision control compares
single runs and no multi-seed study was performed, so the difference between
P12 and P10 cannot be generalised into a stable precision advantage. The 8-D
state is not a network condition in this configuration, which restricts
conclusions about state-conditioned policies in general. Videos are written at
10\,fps for viewing and the last partial action chunk may not be fully
executed, which limits frame-by-frame interpretation.

For a more recent benchmark, \emph{LIBERO-Para} provides 2026 language
paraphrase tests with an official $\pi_{0.5}$ guide and is the natural next
target once the current protocol is stable. \emph{LIBERO-Plus} studies visual,
environmental and language perturbations but needs extra assets and a different
number of evaluation runs; the $\pi_0$ results in its tables must not be read as
$\pi_{0.5}$. \emph{RoboCasa365} is a broader kitchen-manipulation candidate but
requires matching weights, cameras and control adaptation, and LIBERO weights
must not be used to climb its leaderboard. The present result is a trustworthy
baseline reproduction, not a new method, and does not automatically transfer to
human first-person egocentric video.

\subsection{Summary}
\label{sec:pi05-summary}

On a single A100 with the authors' released LIBERO weights and the official
OpenPI implementation, $\pi_{0.5}$ reproduces $1935/2000 = 96.75\,\%$ on the
four LIBERO suites, $0.10$ percentage points below the official $96.85\,\%$.
All 64{,}761 inferences and 320{,}142 executed action steps are finite and
verified. A separate highest-precision control reaches $97.00\,\%$ and is kept
as a control rather than as the main table. Six genuine teaching cases, a
physically replayed drawer failure and two clearly separated classes of anomaly
(non-finite engineering failures versus normal task failures) are documented,
together with all limitations of the claim.
